Buchbergers Algorithm is the algorithm that finds a complete gröbner basis. The Buchberger criterion that was carefully proven in the previous sections of this thesis is fully leveraged here. In doing so, the algorithm itself is rather easy to implement, all the work behind the algorithm is in proving that the criterion is sufficient to generate the Gröbner basis. However, the criterion does not imply that the algorithm will terminate in a finite amount of steps. This is important as we want to ensure that a computer can implement the algorithm successfully. 

\begin{algorithm}
\caption{Buchberger Algorithm}
\begin{algorithmic}[1]
\Require Functions $P = \vectorComponents{p}$ that generate $I$.
\Ensure A Gröbner basis $G = \vectorComponentsX{g}{s}$ for $I$, with $P \subseteq G$

\State $ G \gets P$
\State $\hat{G} \gets \lbrace \rbrace$ 

\While{$ G \neq \hat{G} $}
    \State$\hat{G}  \gets G$

    \For{every pair $\lbrace g_i,g_j \rbrace$, $g_i \neq g_j$ in $\hat{G}$}
        \State $r \gets \sPoly{\spolyFunction{g}{i}{j}}{\hat{G}}$
        \If{$r \neq 0$}
            \State $ G \gets G \cup \lbrace r \rbrace $
        \EndIf
    \EndFor
\EndWhile

\State \Return $ G$
\end{algorithmic}
\end{algorithm}

\begin{proof}
    There are two parts to the proof. First, we must prove that the generated set, $G$, actually is a Gröbner basis, and then we must prove that the algorithm terminates. We can prove that $G$ is a Gröbner basis by first noting that that $G \subseteq I$ at every point of the algorithm. Since we are solely using functions that generate $I$, any remainder from division will be on $I$ as well. So by extending $G$ with the remainder, $G \subseteq I$ holds. The fact that the ideal $G$ is a Gröbner basis comes from theorem \ref{thm:BuchbergersCriterion}. The algorithm terminates when $\sPoly{\spolyFunction{g}{i}{j}}{\hat{G}} = 0$ for all pairs which was the condition for the basis to be a Gröbner basis. 
    
    The remaining part is to prove that the algorithm terminates. Consider the remainder,$r$, that is generated from the division of members on $G$. From lemma \ref{lemma:leadTermMembership}, we can note that the $LT(r)$ is not contained in $\hat{G}$, but is added to $G$. Thus, we can state that the $\ideal{LT(\hat{G})} \subseteq \ideal{LT(G)}$ since $LT(r)$ is solely in $LT(G)$. Since we are dealing with an ascending chain of ideals in $\polyRing$, the ascending chain condition can be used (see chapter 2.5 theorem 7 in \cite{Applied}). By the ascending chain condition, there exists an $N \geq 1$ such that $I_N = I_{N+1} = I_{N+2} = \cdots$. Thus, after N steps of the algorithm, further steps will not add any $r$, and the algorithm will terminate. 
\end{proof}

